% Source: book pp. 133--142 (PDF 147--156); solutions pp. 135--151, 175.
\section{Invariant Subspaces}

\subsection{Eigenvalues}

\begin{definition}[Operator]\label{ax:5.1}
A linear map from a vector space to itself is called an \emph{operator}.
\end{definition}

\marginnote{Recall that we defined the notation $\Lin(V)$ to mean
$\Lin(V,V)$.}%
Suppose $T\in\Lin(V)$. If $m\ge2$ and
\[ V = V_1\oplus\cdots\oplus V_m, \]
where each $V_k$ is a nonzero subspace of $V$, then to understand the behavior
of $T$ we only need to understand the behavior of each $T|_{V_k}$; here
$T|_{V_k}$ denotes the restriction of $T$ to the smaller domain $V_k$. Dealing
with $T|_{V_k}$ should be easier than dealing with $T$ because $V_k$ is a
smaller vector space than $V$.

However, if we intend to apply tools useful in the study of operators (such as
taking powers), then we have a problem: $T|_{V_k}$ may not map $V_k$ into
itself; in other words, $T|_{V_k}$ may not be an operator on $V_k$. Thus we are
led to consider only decompositions of $V$ of the form above in which $T$ maps
each $V_k$ into itself. Hence we now give a name to subspaces of $V$ that get
mapped into themselves by $T$.

\begin{definition}[Invariant subspace]\label{ax:5.2}
Suppose $T\in\Lin(V)$. A subspace $U$ of $V$ is called \emph{invariant} under
$T$ if $Tu\in U$ for every $u\in U$.
\end{definition}

Thus $U$ is invariant under $T$ if $T|_U$ is an operator on $U$.

\begin{example}[Subspace invariant under differentiation operator]\label{ax:5.3}
Suppose that $T\in\Lin(\Poly(\R))$ is defined by $Tp=p'$. Then $\Poly_4(\R)$,
which is a subspace of $\Poly(\R)$, is invariant under $T$ because if
$p\in\Poly(\R)$ has degree at most $4$, then $p'$ also has degree at most $4$.
\end{example}

\begin{example}[Four invariant subspaces, not necessarily all different]\label{ax:5.4}
If $T\in\Lin(V)$, then the following subspaces of $V$ are all invariant under
$T$.
\begin{description}[font=\normalfont,align=right,labelwidth=3.4em,
  labelsep=.8em,leftmargin=4.2em,itemsep=1pt,topsep=3pt]
\item[$\{0\}$] The subspace $\{0\}$ is invariant under $T$ because if
  $u\in\{0\}$, then $u=0$ and hence $Tu=0\in\{0\}$.
\item[$V$] The subspace $V$ is invariant under $T$ because if $u\in V$, then
  $Tu\in V$.
\item[$\nul T$] The subspace $\nul T$ is invariant under $T$ because if
  $u\in\nul T$, then $Tu=0$, and hence $Tu\in\nul T$.
\item[$\range T$] The subspace $\range T$ is invariant under $T$ because if
  $u\in\range T$, then $Tu\in\range T$.
\end{description}
\end{example}

Must an operator $T\in\Lin(V)$ have any invariant subspaces other than $\{0\}$
and $V$? Later we will see that this question has an affirmative answer if $V$
is finite-dimensional and $\dim V>1$ (for $\F=\C$) or $\dim V>2$ (for
$\F=\R$); see \ref{ax:5.19} and Exercise~29 in Section~5B.

The previous example noted that $\nul T$ and $\range T$ are invariant under
$T$. However, these subspaces do not necessarily provide easy answers to the
question above about the existence of invariant subspaces other than $\{0\}$
and $V$, because $\nul T$ may equal $\{0\}$ and $\range T$ may equal $V$ (this
happens when $T$ is invertible).

We will return later to a deeper study of invariant subspaces. Now we turn to
an investigation of the simplest possible nontrivial invariant
subspaces---invariant subspaces of dimension one.

Take any $v\in V$ with $v\ne0$ and let $U$ equal the set of all scalar
multiples of $v$:
\[ U = \{\lambda v : \lambda\in\F\} = \spn(v). \]
Then $U$ is a one-dimensional subspace of $V$ (and every one-dimensional
subspace of $V$ is of this form for an appropriate choice of $v$). If $U$ is
invariant under an operator $T\in\Lin(V)$, then $Tv\in U$, and hence there is
a scalar $\lambda\in\F$ such that
\[ Tv = \lambda v. \]
Conversely, if $Tv=\lambda v$ for some $\lambda\in\F$, then $\spn(v)$ is a
one-dimensional subspace of $V$ invariant under $T$.

The equation $Tv=\lambda v$, which we have just seen is intimately connected
with one-dimensional invariant subspaces, is important enough that the scalars
$\lambda$ and vectors $v$ satisfying it are given special names.

\begin{definition}[Eigenvalue]\label{ax:5.5}
Suppose $T\in\Lin(V)$. A number $\lambda\in\F$ is called an \emph{eigenvalue}
of $T$ if there exists $v\in V$ such that $v\ne0$ and $Tv=\lambda v$.
\end{definition}

\marginnote{The word \textbf{eigenvalue} is half-German, half-English. The
German prefix \textbf{eigen} means ``own'' in the sense of characterizing an
intrinsic property.}%
In the definition above, we require that $v\ne0$ because every scalar
$\lambda\in\F$ satisfies $T0=\lambda0$.

The comments above show that $V$ has a one-dimensional subspace invariant under
$T$ if and only if $T$ has an eigenvalue.

\begin{example}[Eigenvalue]\label{ax:5.6}
Define an operator $T\in\Lin(\F^3)$ by
\[ T(x,y,z) = (7x+3z,\,3x+6y+9z,\,-6y) \]
for $(x,y,z)\in\F^3$. Then $T(3,1,-1)=(18,6,-6)=6(3,1,-1)$. Thus $6$ is an
eigenvalue of $T$.
\end{example}

\marginnote[-4mm]{Reminder: $I\in\Lin(V)$ is the identity operator. Thus $Iv=v$
for all $v\in V$.}%
The equivalences in the next result, along with many deep results in linear
algebra, are valid only in the context of finite-dimensional vector spaces.

\begin{result}[Equivalent conditions to be an eigenvalue]\label{ax:5.7}
Suppose $V$ is finite-dimensional, $T\in\Lin(V)$, and $\lambda\in\F$. Then the
following are equivalent.
\begin{parts}
\item $\lambda$ is an eigenvalue of $T$.
\item $T-\lambda I$ is not injective.
\item $T-\lambda I$ is not surjective.
\item $T-\lambda I$ is not invertible.
\end{parts}
\end{result}

\begin{proof}
Conditions (a) and (b) are equivalent because the equation $Tv=\lambda v$ is
equivalent to the equation $(T-\lambda I)v=0$. Conditions (b), (c), and (d) are
equivalent by \ref{ax:3.65}.
\end{proof}

\begin{definition}[Eigenvector]\label{ax:5.8}
Suppose $T\in\Lin(V)$ and $\lambda\in\F$ is an eigenvalue of $T$. A vector
$v\in V$ is called an \emph{eigenvector} of $T$ corresponding to $\lambda$ if
$v\ne0$ and $Tv=\lambda v$.
\end{definition}

In other words, a nonzero vector $v\in V$ is an eigenvector of an operator
$T\in\Lin(V)$ if and only if $Tv$ is a scalar multiple of $v$. Because
$Tv=\lambda v$ if and only if $(T-\lambda I)v=0$, a vector $v\in V$ with
$v\ne0$ is an eigenvector of $T$ corresponding to $\lambda$ if and only if
$v\in\nul(T-\lambda I)$.

\begin{example}[Eigenvalues and eigenvectors]\label{ax:5.9}
Suppose $T\in\Lin(\F^2)$ is defined by $T(w,z)=(-z,w)$.
\begin{parts}
\item First consider the case $\F=\R$. Then $T$ is a counterclockwise rotation
  by $90^\circ$ about the origin in $\R^2$. An operator has an eigenvalue if and
  only if there exists a nonzero vector in its domain that gets sent by the
  operator to a scalar multiple of itself. A $90^\circ$ counterclockwise
  rotation of a nonzero vector in $\R^2$ cannot equal a scalar multiple of
  itself. Conclusion: if $\F=\R$, then $T$ has no eigenvalues (and thus has no
  eigenvectors).
\item Now consider the case $\F=\C$. To find eigenvalues of $T$, we must find
  the scalars $\lambda$ such that $T(w,z)=\lambda(w,z)$ has some solution other
  than $w=z=0$. The equation $T(w,z)=\lambda(w,z)$ is equivalent to the
  simultaneous equations\refstepcounter{theorem}%
  \[ -z=\lambda w,\qquad w=\lambda z. \tag*{\thetheorem}\label{ax:5.10} \]
  Substituting the value for $w$ given by the second equation into the first
  equation gives
  \[ -z = \lambda^2 z. \]
  Now $z$ cannot equal $0$ [otherwise \ref{ax:5.10} implies that $w=0$; we are
  looking for solutions to \ref{ax:5.10} such that $(w,z)$ is not the $0$
  vector], so the equation above leads to the equation
  \[ -1 = \lambda^2. \]
  The solutions to this equation are $\lambda=i$ and $\lambda=-i$.

  You can verify that $i$ and $-i$ are eigenvalues of $T$. Indeed, the
  eigenvectors corresponding to the eigenvalue $i$ are the vectors of the form
  $(w,-wi)$, with $w\in\C$ and $w\ne0$. Furthermore, the eigenvectors
  corresponding to the eigenvalue $-i$ are the vectors of the form $(w,wi)$,
  with $w\in\C$ and $w\ne0$.
\end{parts}
\end{example}

In the next proof, we again use the equivalence
\[ Tv=\lambda v \iff (T-\lambda I)v=0. \]

\begin{result}[Linearly independent eigenvectors]\label{ax:5.11}
Suppose $T\in\Lin(V)$. Then every list of eigenvectors of $T$ corresponding to
distinct eigenvalues of $T$ is linearly independent.
\end{result}

\begin{proof}
Suppose the desired result is false. Then there exists a smallest positive
integer $m$ such that there exists a linearly dependent list $v_1,\dots,v_m$ of
eigenvectors of $T$ corresponding to distinct eigenvalues
$\lambda_1,\dots,\lambda_m$ of $T$ (note that $m\ge2$ because an eigenvector
is, by definition, nonzero). Thus there exist $a_1,\dots,a_m\in\F$, none of
which are $0$ (because of the minimality of $m$), such that
\[ a_1v_1+\cdots+a_mv_m = 0. \]
Apply $T-\lambda_mI$ to both sides of the equation above, getting
\[ a_1(\lambda_1-\lambda_m)v_1+\cdots+a_{m-1}(\lambda_{m-1}-\lambda_m)v_{m-1}
   = 0. \]
Because the eigenvalues $\lambda_1,\dots,\lambda_m$ are distinct, none of the
coefficients above equal $0$. Thus $v_1,\dots,v_{m-1}$ is a linearly dependent
list of $m-1$ eigenvectors of $T$ corresponding to distinct eigenvalues,
contradicting the minimality of $m$. This contradiction completes the proof.
\end{proof}

The result above leads to a short proof of the result below, which puts an
upper bound on the number of distinct eigenvalues that an operator can have.

\begin{result}[Operator cannot have more eigenvalues than dimension of vector
space]\label{ax:5.12}
Suppose $V$ is finite-dimensional. Then each operator on $V$ has at most
$\dim V$ distinct eigenvalues.
\end{result}

\begin{proof}
Let $T\in\Lin(V)$. Suppose $\lambda_1,\dots,\lambda_m$ are distinct eigenvalues
of $T$. Let $v_1,\dots,v_m$ be corresponding eigenvectors. Then \ref{ax:5.11}
implies that the list $v_1,\dots,v_m$ is linearly independent. Thus
$m\le\dim V$ (see \ref{ax:2.22}), as desired.
\end{proof}

\subsection{Polynomials Applied to Operators}

The main reason that a richer theory exists for operators (which map a vector
space into itself) than for more general linear maps is that operators can be
raised to powers. In this subsection we define that notion and the concept of
applying a polynomial to an operator. This concept will be the key tool that
we use in the next section when we prove that every operator on a nonzero
finite-dimensional complex vector space has an eigenvalue.

If $T$ is an operator, then $TT$ makes sense (see \ref{ax:3.7}) and is also an
operator on the same vector space as $T$. We usually write $T^2$ instead of
$TT$. More generally, we have the following definition of $T^m$.

\begin{notation}[$T^m$]\label{ax:5.13}
Suppose $T\in\Lin(V)$ and $m$ is a positive integer.
\begin{itemize}
\item $T^m\in\Lin(V)$ is defined by $T^m=\ftimes{T\cdots T}{m\text{ times}}$.
\item $T^0$ is defined to be the identity operator $I$ on $V$.
\item If $T$ is invertible with inverse $T^{-1}$, then $T^{-m}\in\Lin(V)$ is
  defined by
  \[ T^{-m} = (T^{-1})^m. \]
\end{itemize}
\end{notation}

You should verify that if $T$ is an operator, then
\[ T^mT^n = T^{m+n} \quad\text{and}\quad (T^m)^n = T^{mn}, \]
where $m$ and $n$ are arbitrary integers if $T$ is invertible and are
nonnegative integers if $T$ is not invertible.

Having defined powers of an operator, we can now define what it means to apply
a polynomial to an operator.

\begin{notation}[$p(T)$]\label{ax:5.14}
Suppose $T\in\Lin(V)$ and $p\in\Poly(\F)$ is a polynomial given by
\[ p(z) = a_0+a_1z+a_2z^2+\cdots+a_mz^m \]
for all $z\in\F$. Then $p(T)$ is the operator on $V$ defined by
\[ p(T) = a_0I+a_1T+a_2T^2+\cdots+a_mT^m. \]
\end{notation}

This is a new use of the symbol $p$ because we are applying $p$ to operators,
not just elements of $\F$. The idea here is that to evaluate $p(T)$, we simply
replace $z$ with $T$ in the expression defining $p$. Note that the constant
term $a_0$ in $p(z)$ becomes the operator $a_0I$ (which is a reasonable choice
because $a_0=a_0z^0$ and thus we should replace $a_0$ with $a_0T^0$, which
equals $a_0I$).

\begin{example}[A polynomial applied to the differentiation
operator]\label{ax:5.15}
Suppose $D\in\Lin(\Poly(\R))$ is the differentiation operator defined by
$Dq=q'$ and $p$ is the polynomial defined by $p(x)=7-3x+5x^2$. Then
$p(D)=7I-3D+5D^2$. Thus
\[ \bigl(p(D)\bigr)q = 7q-3q'+5q'' \]
for every $q\in\Poly(\R)$.
\end{example}

If we fix an operator $T\in\Lin(V)$, then the function from $\Poly(\F)$ to
$\Lin(V)$ given by $p\mapsto p(T)$ is linear, as you should verify.

\begin{definition}[Product of polynomials]\label{ax:5.16}
If $p,q\in\Poly(\F)$, then $pq\in\Poly(\F)$ is the polynomial defined by
\[ (pq)(z) = p(z)q(z) \]
for all $z\in\F$.
\end{definition}

\marginnote[-12mm]{Informal proof: When a product of polynomials is expanded
using the distributive property, it does not matter whether the symbol is $z$
or $T$.}%
The order does not matter in taking products of polynomials of a single
operator, as shown by (b) in the next result.

\begin{result}[Multiplicative properties]\label{ax:5.17}
Suppose $p,q\in\Poly(\F)$ and $T\in\Lin(V)$. Then
\begin{parts}
\item $(pq)(T)=p(T)q(T)$;
\item $p(T)q(T)=q(T)p(T)$.
\end{parts}
\end{result}

\begin{proof}
\begin{parts}
\item Suppose $p(z)=\sum_{j=0}^m a_jz^j$ and $q(z)=\sum_{k=0}^n b_kz^k$ for
  all $z\in\F$. Then
  \[ (pq)(z) = \sum_{j=0}^m\sum_{k=0}^n a_jb_kz^{j+k}. \]
  Thus
  \begin{align*}
  (pq)(T) &= \sum_{j=0}^m\sum_{k=0}^n a_jb_kT^{j+k}\\
          &= \Bigl(\sum_{j=0}^m a_jT^j\Bigr)\Bigl(\sum_{k=0}^n b_kT^k\Bigr)\\
          &= p(T)q(T).
  \end{align*}
\item Using (a) twice, we have $p(T)q(T)=(pq)(T)=(qp)(T)=q(T)p(T)$.
  \qedhere
\end{parts}
\end{proof}

We observed earlier that if $T\in\Lin(V)$, then the subspaces $\nul T$ and
$\range T$ are invariant under $T$ (see \ref{ax:5.4}). Now we show that the
null space and the range of every polynomial of $T$ are also invariant under
$T$.

\begin{result}[Null space and range of $p(T)$ are invariant under
$T$]\label{ax:5.18}
Suppose $T\in\Lin(V)$ and $p\in\Poly(\F)$. Then $\nul p(T)$ and
$\range p(T)$ are invariant under $T$.
\end{result}

\begin{proof}
Suppose $u\in\nul p(T)$. Then $p(T)u=0$. Thus
\[ \bigl(p(T)\bigr)(Tu) = \bigl(p(T)\,T\bigr)(u) = \bigl(T\,p(T)\bigr)(u)
   = T\bigl(p(T)u\bigr) = T(0) = 0. \]
Hence $Tu\in\nul p(T)$. Thus $\nul p(T)$ is invariant under $T$, as desired.

Suppose $u\in\range p(T)$. Then there exists $v\in V$ such that $u=p(T)v$.
Thus
\[ Tu = T\bigl(p(T)v\bigr) = p(T)(Tv). \]
Hence $Tu\in\range p(T)$. Thus $\range p(T)$ is invariant under $T$, as
desired.
\end{proof}

% ------------------------------------------------------------------ exercises
\begin{exc}

\begin{exercise}
Suppose $T\in\Lin(V)$ and $U$ is a subspace of $V$.
\begin{parts}
\item Prove that if $U\subseteq\nul T$, then $U$ is invariant under $T$.
\item Prove that if $\range T\subseteq U$, then $U$ is invariant under $T$.
\end{parts}
\end{exercise}
\begin{solution}
\begin{parts}
\item If $U\subseteq\nul T$, then for any $u\in U$, we have $T(u)=0\in U$, so
  $U$ is invariant under $T$.
\item If $\range T\subseteq U$, then for any $u\in U$,
  $T(u)\in\range T\subseteq U$, so $U$ is invariant under $T$. \qedhere
\end{parts}
\end{solution}

\begin{exercise}
Suppose that $T\in\Lin(V)$ and $V_1,\dots,V_m$ are subspaces of $V$ invariant
under $T$. Prove that $V_1+\cdots+V_m$ is invariant under $T$.
\end{exercise}
\begin{solution}
Let $v\in V_1+\cdots+V_m$. Then $v=v_1+\cdots+v_m$ for some $v_i\in V_i$. Since
each $V_i$ is invariant under $T$, we have $T(v_i)\in V_i$ for each $i$.
Therefore,
\[ T(v) = T(v_1+\cdots+v_m) = T(v_1)+\cdots+T(v_m) \in V_1+\cdots+V_m, \]
so $V_1+\cdots+V_m$ is invariant under $T$.
\end{solution}

\begin{exercise}
Suppose $T\in\Lin(V)$. Prove that the intersection of every collection of
subspaces of $V$ invariant under $T$ is invariant under $T$.
\end{exercise}
\begin{solution}
Let $\{U_i\}_{i\in I}$ be a collection of subspaces of $V$ invariant under $T$,
and let
\[ U = \bigcap_{i\in I} U_i. \]
Let $u\in U$. Then $u\in U_i$ for all $i\in I$. Since each $U_i$ is invariant
under $T$, we have $T(u)\in U_i$ for all $i\in I$. Therefore, $T(u)\in U$, so
$U$ is invariant under $T$.
\end{solution}

\begin{exercise}
Prove or give a counterexample: If $V$ is finite-dimensional and $U$ is a
subspace of $V$ that is invariant under every operator on $V$, then $U=\{0\}$
or $U=V$.
\end{exercise}
\begin{solution}
If $U=\{0\}$, then we are done. Now suppose $U\ne\{0\}$ so that there exists
nonzero $v_1\in U$. Extend $v_1$ to a basis $v_1,\dots,v_m$ of $V$, and let
$v\in V$. Define the linear operator $T\in\Lin(V)$ by $T(v_1)=v$ and
$T(v_i)=0$ for $i=2,\dots,m$. Since $U$ is invariant under $T$, we have
$T(v_1)=v\in U$. But $v$ was an arbitrary element of $V$, so $U=V$. Therefore,
$U=\{0\}$ or $U=V$.
\end{solution}

\begin{exercise}
Suppose $T\in\Lin(\R^2)$ is defined by $T(x,y)=(-3y,x)$. Find the eigenvalues
of $T$.
\end{exercise}
\begin{solution}
If $(x,y)$ is an eigenvector of $T$ with eigenvalue $\lambda$, then we have the
system of equations
\[ \begin{cases} -3y=\lambda x,\\ x=\lambda y \end{cases} \]
Solving the system, we get $\lambda^2=-3$ which has no real solutions.
Therefore, $T$ has no real eigenvalues.
\end{solution}

\begin{exercise}
Define $T\in\Lin(\F^2)$ by $T(w,z)=(z,w)$. Find all eigenvalues and
eigenvectors of $T$.
\end{exercise}
\begin{solution}
If $(w,z)$ is an eigenvector of $T$ with eigenvalue $\lambda$, then we have
\[ (z,w) = \lambda(w,z), \]
which gives the system
\[ \begin{cases} z=\lambda w,\\ w=\lambda z. \end{cases} \]
Solving, we get $\lambda^2=1$, so $\lambda=1$ or $\lambda=-1$. If $\lambda=1$,
then $z=w$, so the eigenvectors are of the form $(w,w)$. If $\lambda=-1$, then
$z=-w$, so the eigenvectors are of the form $(w,-w)$.
\end{solution}

\begin{exercise}
Define $T\in\Lin(\F^3)$ by $T(z_1,z_2,z_3)=(2z_2,0,5z_3)$. Find all eigenvalues
and eigenvectors of $T$.
\end{exercise}
\begin{solution}
If $(z_1,z_2,z_3)$ is an eigenvector of $T$ with eigenvalue $\lambda$, then we
have
\[ (2z_2,0,5z_3) = \lambda(z_1,z_2,z_3), \]
which gives the system
\[ \begin{cases} 2z_2=\lambda z_1,\\ 0=\lambda z_2,\\ 5z_3=\lambda z_3.
   \end{cases} \]
Solving, we get the eigenvalues $\lambda=0$ and $\lambda=5$. For $\lambda=0$,
the eigenvectors are of the form $(z_1,0,0)$. For $\lambda=5$, the eigenvectors
are of the form $(0,0,z_3)$.
\end{solution}

\begin{exercise}
Suppose $P\in\Lin(V)$ is such that $P^2=P$. Prove that if $\lambda$ is an
eigenvalue of $P$, then $\lambda=0$ or $\lambda=1$.
\end{exercise}
\begin{solution}
If $v$ is an eigenvector of $P$ with eigenvalue $\lambda$, then $Pv=\lambda v$.
Applying $P$ again, we get $P^2v=P(\lambda v)=\lambda Pv=\lambda^2v$. But
$P^2=P$, so $P^2v=Pv=\lambda v$. Therefore, $\lambda^2v=\lambda v$, which gives
$(\lambda^2-\lambda)v=0$. Since $v\ne0$, we must have $\lambda^2-\lambda=0$, so
$\lambda(\lambda-1)=0$. Hence, $\lambda=0$ or $\lambda=1$.
\end{solution}

\begin{exercise}
Define $T\colon\Poly(\R)\to\Poly(\R)$ by $Tp=p'$. Find all eigenvalues and
eigenvectors of $T$.
\end{exercise}
\begin{solution}
Let $p\in\Poly(\R)$ such that $p'=\lambda p$ for some $\lambda\in\R$.
\begin{itemize}
\item If $\lambda\ne0$, then $\deg p'=\deg p$, which is a contradiction since
  $\deg p'=\deg p-1$. Therefore, there are no eigenvectors corresponding to
  nonzero eigenvalues.
\item If $\lambda=0$, then $p'=0$, which implies that $p$ is a constant
  polynomial. Therefore, the eigenvectors corresponding to the eigenvalue
  $\lambda=0$ are all nonzero constant polynomials. \qedhere
\end{itemize}
\end{solution}

\begin{exercise}
Define $T\in\Lin(\Poly_4(\R))$ by $(Tp)(x)=xp'(x)$ for all $x\in\R$. Find all
eigenvalues and eigenvectors of $T$.
\end{exercise}
\begin{solution}
Let $p(x)=a_0+a_1x+a_2x^2+a_3x^3+a_4x^4$ be an eigenvector of $T$ with
eigenvalue $\lambda$. Then
\[ (Tp)(x) = xp'(x) = x(a_1+2a_2x+3a_3x^2+4a_4x^3)
   = a_1x+2a_2x^2+3a_3x^3+4a_4x^4. \]
Equating coefficients with
$\lambda p(x)=\lambda a_0+\lambda a_1x+\lambda a_2x^2+\lambda a_3x^3
+\lambda a_4x^4$ yields the system:
\[ \lambda a_0=0,\quad \lambda a_1=a_1,\quad \lambda a_2=2a_2,\quad
   \lambda a_3=3a_3,\quad \lambda a_4=4a_4. \]
For each $n\in\{0,1,2,3,4\}$, the equation $\lambda a_n=na_n$ is satisfied
with $a_n\ne0$ if and only if $\lambda=n$. When $\lambda=n$, all other
coefficients must be zero, so the eigenvector is a nonzero scalar multiple of
$x^n$. Therefore, the eigenvalues of $T$ are $0,1,2,3,4$, with corresponding
eigenvectors being nonzero scalar multiples of $1,x,x^2,x^3,x^4$,
respectively.
\end{solution}

\begin{exercise}
Suppose $V$ is finite-dimensional, $T\in\Lin(V)$, and $\alpha\in\F$. Prove that
there exists $\delta>0$ such that $T-\lambda I$ is invertible for all
$\lambda\in\F$ such that $0<\abs{\alpha-\lambda}<\delta$.
\end{exercise}
\begin{solution}
Since $V$ is finite-dimensional, the set of eigenvalues of $T$ is finite. If
$T$ has no eigenvalues, then $T-\lambda I$ is invertible for all
$\lambda\in\F$, and we can choose any $\delta>0$. Otherwise, let
$\lambda_1,\dots,\lambda_m$ be the eigenvalues of $T$. Define
\[ \delta = \min_{\lambda_j\ne\alpha}\abs{\alpha-\lambda_j}. \]
If $\alpha$ is the only eigenvalue, then we can select any $\delta>0$ since any
$\lambda$ satisfying $0<\abs{\alpha-\lambda}<\delta$ will not be equal to
$\alpha$, and hence $T-\lambda I$ will be invertible. Otherwise, we have for
any $\lambda\in\F$ such that $0<\abs{\alpha-\lambda}<\delta$, we have
$\lambda\ne\lambda_j$ for all $j=1,\dots,m$, so $\lambda$ is not an eigenvalue
of $T$. Hence, $T-\lambda I$ is invertible for all such $\lambda$.
\end{solution}

\begin{exercise}
Suppose $V=U\oplus W$, where $U$ and $W$ are nonzero subspaces of $V$. Define
$P\in\Lin(V)$ by $P(u+w)=u$ for each $u\in U$ and each $w\in W$. Find all
eigenvalues and eigenvectors of $P$.
\end{exercise}
\begin{solution}
Let $v=u+w\in V$ be an eigenvector of $P$ with eigenvalue $\lambda$. Then
$P(v)=P(u+w)=u=\lambda v=\lambda(u+w)$. Comparing the $U$ and $W$ components,
we get $u=\lambda u$ and $0=\lambda w$. If $\lambda\ne0$, then $w=0$ and
$u=\lambda u$ implies $\lambda=1$. If $\lambda=0$, then $u=0$ and $w$ can be
any nonzero vector in $W$. Therefore, the eigenvalues are $\lambda=1$ and
$\lambda=0$, with corresponding eigenvectors being the nonzero vectors in $U$
for $\lambda=1$ and the nonzero vectors in $W$ for $\lambda=0$.
\end{solution}

\begin{exercise}
Suppose $T\in\Lin(V)$. Suppose $S\in\Lin(V)$ is invertible.
\begin{parts}
\item Prove that $T$ and $S^{-1}TS$ have the same eigenvalues.
\item What is the relationship between the eigenvectors of $T$ and the
  eigenvectors of $S^{-1}TS$?
\end{parts}
\end{exercise}
\begin{solution}
\begin{parts}
\item Let $\lambda$ be an eigenvalue of $T$ with eigenvector $v\in V$, so
  $T(v)=\lambda v$. Then
  \[ (S^{-1}TS)(S^{-1}v) = S^{-1}T(v) = S^{-1}(\lambda v)
     = \lambda(S^{-1}v). \]
  Thus, $S^{-1}v$ is an eigenvector of $S^{-1}TS$ with eigenvalue $\lambda$.
  This shows that every eigenvalue of $T$ is an eigenvalue of $S^{-1}TS$.
  Conversely, if $\lambda$ is an eigenvalue of $S^{-1}TS$ with eigenvector $u$,
  then $Su$ is an eigenvector of $T$ with eigenvalue $\lambda$. Therefore, $T$
  and $S^{-1}TS$ have the same eigenvalues.
\item The relation between the eigenvectors of $T$ and the eigenvectors of
  $S^{-1}TS$ is given by $v\mapsto S^{-1}v$. Specifically, if $v$ is an
  eigenvector of $T$ with eigenvalue $\lambda$, then $S^{-1}v$ is an
  eigenvector of $S^{-1}TS$ with the same eigenvalue $\lambda$. Conversely, if
  $u$ is an eigenvector of $S^{-1}TS$ with eigenvalue $\lambda$, then $Su$ is
  an eigenvector of $T$ with eigenvalue $\lambda$. \qedhere
\end{parts}
\end{solution}

\begin{exercise}
Give an example of an operator on $\R^4$ that has no (real) eigenvalues.
\end{exercise}
\begin{solution}
Consider the linear operator on $\R^4$ given by the matrix
\[ T = \begin{pmatrix} 0&-1&0&0\\ 1&0&0&0\\ 0&0&0&-1\\ 0&0&1&0 \end{pmatrix}.
\]
Letting $T$ act on $(x_1,x_2,x_3,x_4)\in\R^4$, we have
\[ T(x_1,x_2,x_3,x_4) = (-x_2,x_1,-x_4,x_3). \]
Then any eigenvector $(x_1,x_2,x_3,x_4)\in\R^4$ with eigenvalue $\lambda$ must
satisfy
\[ (-x_2,x_1,-x_4,x_3) = \lambda(x_1,x_2,x_3,x_4). \]
This yields $\lambda^2+1=0$, which has no solution in $\R$.
\end{solution}

\begin{exercise}
Suppose $V$ is finite-dimensional, $T\in\Lin(V)$, and $\lambda\in\F$. Show that
$\lambda$ is an eigenvalue of $T$ if and only if $\lambda$ is an eigenvalue of
the dual operator $T'\in\Lin(V')$.
\end{exercise}
\begin{solution}
\begin{align*}
\lambda\text{ is an eigenvalue of }T
  &\iff T-\lambda I\text{ is not invertible}\\
  &\iff (T-\lambda I)'\text{ is not invertible}\\
  &\iff T'-\lambda I\text{ is not invertible}\\
  &\iff \lambda\text{ is an eigenvalue of }T'. \qedhere
\end{align*}
\end{solution}

\begin{exercise}
Suppose $v_1,\dots,v_n$ is a basis of $V$ and $T\in\Lin(V)$. Prove that if
$\lambda$ is an eigenvalue of $T$, then
\[ \abs{\lambda} \le n\max\bigl\{\abs{\Mat(T)_{j,k}} : 1\le j,k\le n\bigr\}, \]
where $\Mat(T)_{j,k}$ denotes the entry in row $j$, column $k$ of the matrix of
$T$ with respect to the basis $v_1,\dots,v_n$.
\exnote{See Exercise~19 in Section~6A for a different bound on
$\abs{\lambda}$.}
\end{exercise}
\begin{solution}
Let $\lambda$ be an eigenvalue of $T$ with eigenvector $u\in V$, so
$Tu=\lambda u$. Let $\Mat(T)$ be the matrix of $T$ with respect to the basis
$v_1,\dots,v_n$. Then
\[ \lambda u_j = \bigl(\Mat(T)u\bigr)_j = \sum_{k=1}^n\Mat(T)_{j,k}u_k
   \quad\text{for each } j=1,\dots,n. \]
Taking absolute values, we get
\begin{align*}
\abs{\lambda}\abs{u_j}
  &= \biggl\lvert\sum_{k=1}^n\Mat(T)_{j,k}u_k\biggr\rvert
   \le \sum_{k=1}^n\abs{\Mat(T)_{j,k}}\abs{u_k}\\
  &\le n\max\bigl\{\abs{\Mat(T)_{j,k}} : 1\le j,k\le n\bigr\}
   \max\bigl\{\abs{u_k} : 1\le k\le n\bigr\}.
\end{align*}
Choosing $j$ such that $\abs{u_j}=\max\{\abs{u_k} : 1\le k\le n\}$, we obtain
\[ \abs{\lambda} \le n\max\bigl\{\abs{\Mat(T)_{j,k}} : 1\le j,k\le n\bigr\}.
   \qedhere \]
\end{solution}

\begin{exercise}
Suppose $\F=\R$, $T\in\Lin(V)$, and $\lambda\in\R$. Prove that $\lambda$ is an
eigenvalue of $T$ if and only if $\lambda$ is an eigenvalue of the
complexification $T_{\C}$.
\exnote{See Exercise~33 in Section~3B for the definition of $T_{\C}$.}
\end{exercise}
\begin{solution}
Suppose $\lambda$ is an eigenvalue of $T$ with eigenvector $v\in V$, so
$Tv=\lambda v$. Then in the complexification, we have
$T_{\C}(v+i0)=Tv+iT0=\lambda v+i0=\lambda(v+i0)$, showing that $\lambda$ is an
eigenvalue of $T_{\C}$ with eigenvector $v+i0$.

Conversely, suppose $\lambda$ is an eigenvalue of $T_{\C}$ with eigenvector
$v+iw\in V_{\C}$, where $v,w\in V$. Then
\[ T_{\C}(v+iw) = Tv+iTw = \lambda(v+iw) = \lambda v+i\lambda w. \]
Equating real and imaginary parts, we get
\[ Tv=\lambda v \quad\text{and}\quad Tw=\lambda w. \]
Since $\lambda\in\R$, if $v\ne0$, then $v$ is an eigenvector of $T$
corresponding to $\lambda$. If $v=0$, then $w\ne0$ and $Tw=\lambda w$, so $w$
is an eigenvector of $T$ corresponding to $\lambda$. This shows that $\lambda$
is an eigenvalue of $T$.
\end{solution}

\begin{exercise}
Suppose $\F=\R$, $T\in\Lin(V)$, and $\lambda\in\C$. Prove that $\lambda$ is an
eigenvalue of the complexification $T_{\C}$ if and only if $\overline{\lambda}$
is an eigenvalue of $T_{\C}$.
\end{exercise}
\begin{solution}
Suppose $\lambda$ is an eigenvalue of $T_{\C}$ with eigenvector
$v+iw\in V_{\C}$, where $v,w\in V$. Then
\[ T_{\C}(v+iw) = Tv+iTw = \lambda(v+iw) = \lambda v+i\lambda w. \]
Moreover, we have
\[ T_{\C}(v-iw) = Tv-iTw. \]
Note that $v+iw=0$ if and only if both $v=w=0$. Then $v+iw\ne0$ implies
$v-iw\ne0$. Let $\lambda=a+bi$ with $a,b\in\R$. Then equating real and
imaginary parts, we have
\[ Tv=av-bw \quad\text{and}\quad Tw=bv+aw. \]
Then
\begin{align*}
T_{\C}(v-iw) &= Tv-iTw\\
  &= (av-bw)-i(bv+aw)\\
  &= av-ibv-bw-iaw\\
  &= (a-ib)(v-iw)\\
  &= \overline{\lambda}(v-iw).
\end{align*}
Since $v-iw\ne0$, this shows that $\overline{\lambda}$ is an eigenvalue of
$T_{\C}$ with eigenvector $v-iw$.

Conversely, if $\overline{\lambda}$ is an eigenvalue of $T_{\C}$, then by the
same argument, $\lambda$ is an eigenvalue of $T_{\C}$. This completes the
proof.
\end{solution}

\begin{exercise}
Show that the forward shift operator $T\in\Lin(\F^\infty)$ defined by
\[ T(z_1,z_2,\dots) = (0,z_1,z_2,\dots) \]
has no eigenvalues.
\end{exercise}
\begin{solution}
Suppose $\lambda\in\F$ is an eigenvalue of $T$ with eigenvector
$(z_1,z_2,\dots)\in\F^\infty$. Then
\[ T(z_1,z_2,\dots) = (0,z_1,z_2,\dots) = \lambda(z_1,z_2,\dots). \]
Equating the first component, we get $0=\lambda z_1$, so either $\lambda=0$ or
$z_1=0$. If $\lambda=0$, then the second component gives $z_1=0$, and the third
component gives $z_2=0$, and so on. By induction, all $z_k=0$, contradicting
that the eigenvector is nonzero. If $\lambda\ne0$, then $z_1=0$, and the second
component gives $z_2=0$, and so on. Again, all $z_k=0$, a contradiction.
Therefore, $T$ has no eigenvalues.
\end{solution}

\begin{exercise}
Define the backward shift operator $S\in\Lin(\F^\infty)$ by
\[ S(z_1,z_2,z_3,\dots) = (z_2,z_3,\dots). \]
\begin{parts}
\item Show that every element of $\F$ is an eigenvalue of $S$.
\item Find all eigenvectors of $S$.
\end{parts}
\end{exercise}
\begin{solution}
\begin{parts}
\item Let $\alpha\in\F$. Consider the vector $(1,\alpha,\alpha^2,\dots)\in
  \F^\infty$. Then
  \[ S(1,\alpha,\alpha^2,\dots) = (\alpha,\alpha^2,\alpha^3,\dots)
     = \alpha(1,\alpha,\alpha^2,\dots), \]
  showing that $\alpha$ is an eigenvalue of $S$ with eigenvector
  $(1,\alpha,\alpha^2,\dots)$. Since $\alpha\in\F$ was arbitrary, every
  element of $\F$ is an eigenvalue of $S$.
\item Let $\alpha\in\F$ and consider the eigenvalue $\alpha$ of $S$. Any
  eigenvector corresponding to $\alpha$ must satisfy
  \[ S(z_1,z_2,z_3,\dots) = \alpha(z_1,z_2,z_3,\dots), \]
  which gives
  \[ (z_2,z_3,z_4,\dots) = (\alpha z_1,\alpha z_2,\alpha z_3,\dots). \]
  Equating components, we get $z_2=\alpha z_1$,
  $z_3=\alpha z_2=\alpha^2z_1$, and in general $z_k=\alpha^{k-1}z_1$ for
  $k\ge1$. Therefore, all eigenvectors corresponding to $\alpha$ are of the
  form
  $(z_1,\alpha z_1,\alpha^2z_1,\dots)=z_1(1,\alpha,\alpha^2,\dots)$ with
  $z_1\ne0$. \qedhere
\end{parts}
\end{solution}

\begin{exercise}
Suppose $T\in\Lin(V)$ is invertible.
\begin{parts}
\item Suppose $\lambda\in\F$ with $\lambda\ne0$. Prove that $\lambda$ is an
  eigenvalue of $T$ if and only if $1/\lambda$ is an eigenvalue of
  $T^{-1}$.
\item Prove that $T$ and $T^{-1}$ have the same eigenvectors.
\end{parts}
\end{exercise}
\begin{solution}
\begin{parts}
\item Suppose $\lambda\in\F$ with $\lambda\ne0$ is an eigenvalue of $T$ with
  eigenvector $v\in V$. Then
  \[ Tv = \lambda v. \]
  Since $T$ is invertible, we can apply $T^{-1}$ to both sides to get
  \[ v = T^{-1}(\lambda v) = \lambda T^{-1}v. \]
  Dividing both sides by $\lambda$ (which is nonzero), we obtain
  \[ T^{-1}v = \frac{1}{\lambda}v. \]
  This shows that $1/\lambda$ is an eigenvalue of $T^{-1}$ with the
  same eigenvector $v$.

  Conversely, suppose $\mu\in\F$ with $\mu\ne0$ is an eigenvalue of $T^{-1}$
  with eigenvector $v\in V$. Then
  \[ T^{-1}v = \mu v. \]
  Applying $T$ to both sides gives
  \[ v = \mu Tv \implies Tv = \frac{1}{\mu}v. \]
  This shows that $1/\mu$ is an eigenvalue of $T$ with the same
  eigenvector $v$.

  Therefore, $\lambda$ is an eigenvalue of $T$ if and only if
  $1/\lambda$ is an eigenvalue of $T^{-1}$.
\item Part (a) shows that if $\lambda$ is an eigenvalue of $T$, then
  $1/\lambda$ is an eigenvalue of $T^{-1}$ with the same eigenvector.
  Conversely, if $\mu$ is an eigenvalue of $T^{-1}$, then $1/\mu$ is an
  eigenvalue of $T$ with the same eigenvector. Therefore, $T$ and $T^{-1}$ have
  the same eigenvectors. \qedhere
\end{parts}
\end{solution}

\begin{exercise}
Suppose $T\in\Lin(V)$ and there exist nonzero vectors $u$ and $w$ in $V$ such
that
\[ Tu=3w \quad\text{and}\quad Tw=3u. \]
Prove that $3$ or $-3$ is an eigenvalue of $T$.
\end{exercise}
\begin{solution}
Consider the vectors $u+w$ and $u-w$. Since $u$ and $w$ are both nonzero, at
least one of $u+w$ or $u-w$ is nonzero. We have
\[ T(u+w) = Tu+Tw = 3w+3u = 3(u+w), \]
and
\[ T(u-w) = Tu-Tw = 3w-3u = -3(u-w). \]
Hence, one of $u+w$ or $u-w$ is an eigenvector with eigenvalue $3$ or $-3$,
respectively.
\end{solution}

\begin{exercise}
Suppose $V$ is finite-dimensional and $S,T\in\Lin(V)$. Prove that $ST$ and
$TS$ have the same eigenvalues.
\end{exercise}
\begin{solution}
Let $\lambda$ be an eigenvalue of $ST$ with eigenvector $v\in V$, so that
\[ STv = \lambda v. \]
Let $w=Tv$. Then
\[ TSw = T(STv) = T(\lambda v) = \lambda Tv = \lambda w. \]
If $\lambda w=0$, then either $\lambda=0$ or $w=0$. If $w=0$, then $Tv=0$,
which implies $STv=S(Tv)=S0=0$. But $STv=\lambda v$, so $\lambda v=0$. Since
$v\ne0$, we must have $\lambda=0$. Then this implies $ST$ is not invertible so
that $TS$ is also not invertible, meaning $0$ is an eigenvalue of $TS$.
Therefore, every eigenvalue of $ST$ is an eigenvalue of $TS$. By symmetry, the
converse also holds, so $ST$ and $TS$ have the same eigenvalues.
\end{solution}

\begin{exercise}
Suppose $A$ is an $n$-by-$n$ matrix with entries in $\F$. Define
$T\in\Lin(\F^n)$ by $Tx=Ax$, where elements of $\F^n$ are thought of as
$n$-by-$1$ column vectors.
\begin{parts}
\item Suppose the sum of the entries in each row of $A$ equals $1$. Prove that
  $1$ is an eigenvalue of $T$.
\item Suppose the sum of the entries in each column of $A$ equals $1$. Prove
  that $1$ is an eigenvalue of $T$.
\end{parts}
\end{exercise}
\begin{solution}
\begin{parts}
\item Suppose the sum of the entries in each row of $A$ equals $1$. Let
  $v=(1,1,\dots,1)^t\in\F^n$. Then
  \[ Tv = Av = v, \]
  showing that $v$ is an eigenvector of $T$ corresponding to the eigenvalue
  $1$.
\item Suppose the sum of the entries in each column of $A$ equals $1$. By part
  (a), this is equivalent to saying that the sum of the entries in each row of
  $A^t$ equals $1$, hence $1$ is an eigenvalue of the linear operator defined
  by $A^t$. But the eigenvalues of $A^t$ are the same as the eigenvalues of $A$
  because $A-\lambda I$ is not invertible if and only if $A^t-\lambda I$ is not
  invertible. Therefore, $1$ is an eigenvalue of $T$. \qedhere
\end{parts}
\end{solution}

\begin{exercise}
Suppose $T\in\Lin(V)$ and $u,w$ are eigenvectors of $T$ such that $u+w$ is
also an eigenvector of $T$. Prove that $u$ and $w$ are eigenvectors of $T$
corresponding to the same eigenvalue.
\end{exercise}
\begin{solution}
Let $Tu=\lambda u$ and $Tw=\mu w$ for some scalars $\lambda,\mu\in\F$. Since
$u+w$ is also an eigenvector of $T$, we have $T(u+w)=\nu(u+w)$ for some scalar
$\nu\in\F$. Then
\[ \lambda u+\mu w = Tu+Tw = T(u+w) = \nu(u+w) = \nu u+\nu w. \]
Since $u$ and $w$ are eigenvectors, they are nonzero. If $u$ and $w$ were
linearly dependent, let $u=\alpha w$ for some scalar $\alpha\in\F$. Then
$Tu=\lambda u=\lambda\alpha w$, and $Tu=T(\alpha w)=\alpha Tw=\alpha\mu w$.
Equating these, we get $\lambda\alpha w=\alpha\mu w$, and since $\alpha\ne0$
and $w\ne0$, it follows that $\lambda=\mu$. Hence, $u$ and $w$ correspond to
the same eigenvalue.

If $u$ and $w$ are linearly independent, we can equate coefficients of $u$ and
$w$ in $\lambda u+\mu w=\nu u+\nu w$ to get $\lambda=\nu=\mu$. Therefore, in
either case, $u$ and $w$ are eigenvectors corresponding to the same
eigenvalue.
\end{solution}

\begin{exercise}
Suppose $T\in\Lin(V)$ is such that every nonzero vector in $V$ is an
eigenvector of $T$. Prove that $T$ is a scalar multiple of the identity
operator.
\end{exercise}
\begin{solution}
Let $v\in V$ be any nonzero vector. By assumption, $v$ is an eigenvector of
$T$, so $Tv=\lambda_vv$ for some scalar $\lambda_v\in\F$. Now, let $u\in V$ be
another nonzero vector. Consider the vector $u+v$. If $u$ and $v$ are linearly
dependent, then without loss of generality, let $u=\alpha v$ for some scalar
$\alpha\in\F$. Then $Tu=\lambda_uu=\lambda_u\alpha v$, and
$Tu=T(\alpha v)=\alpha Tv=\alpha\lambda_vv$. Equating these, we get
$\lambda_u\alpha v=\alpha\lambda_vv$, and since $\alpha\ne0$ and $v\ne0$, it
follows that $\lambda_u=\lambda_v$.

The remaining case is when $u+v$ is nonzero such that they are linearly
independent. In this case, $u+v$ is an eigenvector of $T$, so
$T(u+v)=\lambda_{u+v}(u+v)$ for some scalar $\lambda_{u+v}\in\F$. On the other
hand, $T(u+v)=Tu+Tv=\lambda_uu+\lambda_vv$. Equating these, we get
$\lambda_{u+v}u+\lambda_{u+v}v=\lambda_uu+\lambda_vv$. Since $u$ and $v$ are
linearly independent, we can equate coefficients to get
$\lambda_{u+v}=\lambda_u=\lambda_v$.

Since $u$ and $v$ were arbitrary nonzero vectors, it follows that there exists
a scalar $\lambda\in\F$ such that $Tv=\lambda v$ for all nonzero $v\in V$.
Hence, $T=\lambda I$.
\end{solution}

\begin{exercise}
Suppose that $V$ is finite-dimensional and $k\in\{1,\dots,\dim V-1\}$. Suppose
$T\in\Lin(V)$ is such that every subspace of $V$ of dimension $k$ is invariant
under $T$. Prove that $T$ is a scalar multiple of the identity operator.
\end{exercise}
\begin{solution}
Let $v_1\in V$ be nonzero. Our goal is to intersect enough $k$-dimensional
subspaces to isolate $\spn v_1$. We proceed as follows. Let $\dim V=n$, and
extend $v_1$ to a basis $v_1,\dots,v_n$ of $V$. For each $i=2,\dots,n$, let
\[ U_i = \spn(v_1,w_2,\dots,w_k),\quad
   w_2,\dots,w_k\in\{v_2,\dots,v_n\}\setminus\{v_i\}, \]
i.e., each $U_i$ must contain $v_1$ as well as $k-1$ vectors from the basis
except $v_i$. This choice is possible, since $k\le n-1$, so there are enough
vectors in $\{v_2,\dots,v_n\}\setminus\{v_i\}$ to choose $k-1$ of them. The
intersection of all such $U_i$ will be exactly $\spn v_1$. Since each $U_i$ is
invariant under $T$, we have $Tv_1\in U_i$ for all $i$. Therefore,
$Tv_1\in\spn v_1$, which means $Tv_1$ is a scalar multiple of $v_1$. Since
$v_1$ was arbitrary, it follows that $T$ is a scalar multiple of the identity
operator.
\end{solution}

\begin{exercise}
Suppose $V$ is finite-dimensional and $T\in\Lin(V)$. Prove that $T$ has at
most $1+\dim\range T$ distinct eigenvalues.
\end{exercise}
\begin{solution}
Let $\lambda_1,\dots,\lambda_m$ be distinct eigenvalues of $T$ with
corresponding eigenvectors $v_1,\dots,v_m$. Since the eigenvalues are
distinct, at most one eigenvalue must be $0$. Without loss of generality,
assume $\lambda_1=0$ if $0$ is an eigenvalue. Then $v_2,\dots,v_m$ are $m-1$
eigenvectors corresponding to nonzero eigenvalues. Since
$Tv_i=\lambda_iv_i\ne0$ for $i=2,\dots,m$, it follows that $v_i\in\range T$.
Therefore, $v_2,\dots,v_m$ are linearly independent vectors in $\range T$,
which implies that $m-1\le\dim\range T$. Hence, the total number of distinct
eigenvalues is at most $1+\dim\range T$.
\end{solution}

\begin{exercise}
Suppose $T\in\Lin(\R^3)$ and $-4$, $5$, and $\sqrt7$ are eigenvalues of $T$.
Prove that there exists $x\in\R^3$ such that $Tx-9x=(-4,5,\sqrt7)$.
\end{exercise}
\begin{solution}
Since $\dim\R^3=3$, the eigenvalues of $T$ are exactly $-4$, $5$, and
$\sqrt7$. It follows that $9$ is not an eigenvalue of $T$, and
$\nul(T-9I)=\{0\}$. Then $T-9I$ is injective, and since $\R^3$ is
finite-dimensional, $T-9I$ is also surjective. Therefore, for any vector in
$\R^3$, there exists $x\in\R^3$ such that $(T-9I)x=(-4,5,\sqrt7)$.
\end{solution}

\begin{exercise}
Suppose $T\in\Lin(V)$ and $(T-2I)(T-3I)(T-4I)=0$. Suppose $\lambda$ is an
eigenvalue of $T$. Prove that $\lambda=2$ or $\lambda=3$ or $\lambda=4$.
\end{exercise}
\begin{solution}
Let $v$ be an eigenvector of $T$ corresponding to the eigenvalue $\lambda$.
Then
\begin{align*}
(T-2I)(T-3I)(T-4I)v &= (T-2I)(T-3I)(\lambda-4)v\\
  &= (T-2I)(\lambda-3)(\lambda-4)v\\
  &= (\lambda-2)(\lambda-3)(\lambda-4)v = 0.
\end{align*}
Since $v\ne0$, it must be that $(\lambda-2)(\lambda-3)(\lambda-4)=0$.
Therefore, $\lambda=2$, $\lambda=3$, or $\lambda=4$.
\end{solution}

\begin{exercise}
Give an example of $T\in\Lin(\R^2)$ such that $T^4=-I$.
\end{exercise}
\begin{solution}
Consider the rotation matrix by $\pi/4$ radians:
\[ T = \begin{pmatrix} \cos(\pi/4) & -\sin(\pi/4)\\
   \sin(\pi/4) & \cos(\pi/4) \end{pmatrix}. \]
Then
\[ T^4 = \begin{pmatrix} \cos(\pi) & -\sin(\pi)\\ \sin(\pi) & \cos(\pi)
   \end{pmatrix} = \begin{pmatrix} -1&0\\ 0&-1 \end{pmatrix} = -I. \qedhere \]
\end{solution}

\begin{exercise}
Suppose $T\in\Lin(V)$ has no eigenvalues and $T^4=I$. Prove that $T^2=-I$.
\end{exercise}
\begin{solution}
\solnote{Manual Remark}{The proof of this is a little easier if we may use
minimal polynomials.}
Note that
\[ T^4-I = (T^2-I)(T^2+I) = (T-I)(T+I)(T^2+I) = 0. \]
Applying $T^4-I$ to any nonzero $v\in V$, we get
\[ (T-I)(T+I)(T^2+I)v = 0. \]
Let $(T^2+I)v=w$. If $w\ne0$, then $(T-I)(T+I)w=0$. If $(T+I)w\ne0$, then
\[ (T-I)(T+I)w = T(T+I)w-I(T+I)w = 0, \]
which implies $T(T+I)w=(T+I)w$. Then $(T+I)w$ would have to be an eigenvector
of $T$ corresponding to the eigenvalue $1$, contradicting the assumption that
$T$ has no eigenvalues. If $(T+I)w=0$, then $Tw=-w$, which would make $w$ an
eigenvector of $T$ corresponding to the eigenvalue $-1$, again contradicting
the assumption that $T$ has no eigenvalues. Hence, we must have $w=0$, i.e.,
$(T^2+I)v=0$. Since $v$ was arbitrary and nonzero, we conclude that
$T^2+I=0$, i.e., $T^2=-I$.
\end{solution}

\begin{exercise}
Suppose $T\in\Lin(V)$ and $m$ is a positive integer.
\begin{parts}
\item Prove that $T$ is injective if and only if $T^m$ is injective.
\item Prove that $T$ is surjective if and only if $T^m$ is surjective.
\end{parts}
\end{exercise}
\begin{solution}
\begin{parts}
\item Suppose $T$ is injective. Then if $T^mv=0$, we have $T(T^{m-1}v)=0$,
  which implies $T^{m-1}v=0$ since $T$ is injective. Repeating this argument,
  we eventually get $v=0$. Hence, $T^m$ is injective.

  Conversely, suppose $T^m$ is injective. If $Tv=0$, then $T^mv=0$, which
  implies $v=0$. Hence, $T$ is injective.
\item Suppose $T$ is surjective. Then for any $v\in V$, there exists $w\in V$
  such that $Tw=v$. Similarly, there exists $w_2\in V$ such that $Tw_2=w$, and
  so on. Repeating this $m$ times, we find $u\in V$ such that $T^mu=v$. Hence,
  $T^m$ is surjective.

  Conversely, suppose $T^m$ is surjective. Then for any $v\in V$, there exists
  $u\in V$ such that $T^mu=v$. Let $w=T^{m-1}u$. Then
  $Tw=T(T^{m-1}u)=T^mu=v$. Hence, $T$ is surjective. \qedhere
\end{parts}
\end{solution}

\begin{exercise}
Suppose $V$ is finite-dimensional and $v_1,\dots,v_m\in V$. Prove that the list
$v_1,\dots,v_m$ is linearly independent if and only if there exists
$T\in\Lin(V)$ such that $v_1,\dots,v_m$ are eigenvectors of $T$ corresponding
to distinct eigenvalues.
\end{exercise}
\begin{solution}
Suppose $v_1,\dots,v_m$ is linearly independent. Since $V$ is
finite-dimensional, we can extend this list to a basis
$v_1,\dots,v_m,v_{m+1},\dots,v_n$ of $V$. Define $T\in\Lin(V)$ by setting
$Tv_i=\lambda_iv_i$ for $i=1,\dots,m$, where $\lambda_1,\dots,\lambda_m$ are
distinct scalars, and choose $\lambda_{m+1},\dots,\lambda_n$ arbitrarily
(distinct from each other and from $\lambda_1,\dots,\lambda_m$). Then
$v_1,\dots,v_m$ are eigenvectors of $T$ corresponding to distinct
eigenvalues.

Conversely, suppose there exists $T\in\Lin(V)$ such that $v_1,\dots,v_m$ are
eigenvectors corresponding to distinct eigenvalues
$\lambda_1,\dots,\lambda_m$. By \ref{ax:5.11}, eigenvectors corresponding to
distinct eigenvalues are linearly independent. Hence, $v_1,\dots,v_m$ are
linearly independent.
\end{solution}

\begin{exercise}
Suppose that $\lambda_1,\dots,\lambda_n$ is a list of distinct real numbers.
Prove that the list $e^{\lambda_1x},\dots,e^{\lambda_nx}$ is linearly
independent in the vector space of real-valued functions on $\R$.
\exnote{Hint: Let $V=\spn(e^{\lambda_1x},\dots,e^{\lambda_nx})$, and define an
operator $D\in\Lin(V)$ by $Df=f'$. Find eigenvalues and eigenvectors of $D$.}
\end{exercise}
\begin{solution}
Suppose there exists a linear combination
\[ a_1e^{\lambda_1x}+\cdots+a_ne^{\lambda_nx} = 0 \]
for all $x\in\R$. Let $V=\spn(e^{\lambda_1x},\dots,e^{\lambda_nx})$ and define
$D\in\Lin(V)$ by $Df=f'$. Then $De^{\lambda_ix}=\lambda_ie^{\lambda_ix}$, so
$e^{\lambda_ix}$ is an eigenvector of $D$ corresponding to the eigenvalue
$\lambda_i$. Since the $\lambda_i$ are distinct, by \ref{ax:5.11}, the
eigenvectors corresponding to distinct eigenvalues are linearly independent.
Hence, $a_1=\cdots=a_n=0$, proving that $e^{\lambda_1x},\dots,e^{\lambda_nx}$
are linearly independent.
\end{solution}

\begin{exercise}
Suppose that $\lambda_1,\dots,\lambda_n$ is a list of distinct positive
numbers. Prove that the list $\cos(\lambda_1x),\dots,\cos(\lambda_nx)$ is
linearly independent in the vector space of real-valued functions on $\R$.
\end{exercise}
\begin{solution}
Suppose there exists a linear combination
\[ a_1\cos(\lambda_1x)+\cdots+a_n\cos(\lambda_nx) = 0 \]
for all $x\in\R$. Let
$V=\spn\bigl(\cos(\lambda_1x),\sin(\lambda_1x),\dots,\cos(\lambda_nx),
\sin(\lambda_nx)\bigr)$ and define $D\in\Lin(V)$ by $Df=f'$. Then
$D^2\cos(\lambda_ix)=-\lambda_i^2\cos(\lambda_ix)$, so $\cos(\lambda_ix)$ is an
eigenvector of $D^2$ corresponding to the eigenvalue $-\lambda_i^2$. Since the
$\lambda_i$ are distinct and positive, the $-\lambda_i^2$ are distinct, and by
\ref{ax:5.11}, the eigenvectors corresponding to distinct eigenvalues are
linearly independent. Hence, $a_1=\cdots=a_n=0$, proving that
$\cos(\lambda_1x),\dots,\cos(\lambda_nx)$ are linearly independent.
\end{solution}

\begin{exercise}
Suppose $V$ is finite-dimensional and $T\in\Lin(V)$. Define
$\mathcal{A}\in\Lin(\Lin(V))$ by
\[ \mathcal{A}(S) = TS \]
for each $S\in\Lin(V)$. Prove that the set of eigenvalues of $T$ equals the
set of eigenvalues of $\mathcal{A}$.
\end{exercise}
\begin{solution}
Suppose $\lambda$ is an eigenvalue of $T$ with corresponding eigenvector
$v\in V$, so $Tv=\lambda v$. Consider a nonzero $\varphi\in V'$ and consider
$S\in\Lin(V)$ defined by $Sw=\varphi(w)v$ for all $w\in V$. Then $S\ne0$ and
for all $w\in V$,
\[ TSw = T\bigl(\varphi(w)v\bigr) = \varphi(w)Tv = \varphi(w)\lambda v
   = \lambda\varphi(w)v = \lambda Sw. \]
Hence, $\lambda$ is an eigenvalue of $\mathcal{A}$.

Conversely, if $\lambda$ is an eigenvalue of $\mathcal{A}$ with corresponding
eigenvector $S\in\Lin(V)$, then $\mathcal{A}(S)=TS=\lambda S$. Since $S\ne0$,
there exists $v\in V$ such that $Sv\ne0$. Then $TSv=\lambda Sv$, so $\lambda$
is an eigenvalue of $T$. Therefore, the set of eigenvalues of $T$ equals the
set of eigenvalues of $\mathcal{A}$.
\end{solution}

\begin{exercise}
Suppose $V$ is finite-dimensional, $T\in\Lin(V)$, and $U$ is a subspace of $V$
invariant under $T$. The \emph{quotient operator} $T/U\in\Lin(V/U)$ is defined
by
\[ (T/U)(v+U) = Tv+U \]
for each $v\in V$.
\begin{parts}
\item Show that the definition of $T/U$ makes sense (which requires using the
  condition that $U$ is invariant under $T$) and show that $T/U$ is an
  operator on $V/U$.
\item Show that each eigenvalue of $T/U$ is an eigenvalue of $T$.
\end{parts}
\end{exercise}
\begin{solution}
\begin{parts}
\item To show that the definition of $T/U$ makes sense, suppose
  $v+U=v'+U$ for some $v,v'\in V$. Then $v-v'\in U$, and since $U$ is
  invariant under $T$, we have $T(v-v')\in U$, so $Tv-Tv'\in U$, which implies
  $Tv+U=Tv'+U$. Hence, the definition of $T/U$ is well-defined.

  To show that $T/U$ is an operator on $V/U$, note that for any
  $v+U\in V/U$, $(T/U)(v+U)=Tv+U\in V/U$. Linearity follows from the linearity
  of $T$: for any $\alpha,\beta\in\F$ and $v+U,w+U\in V/U$,
  \begin{align*}
  (T/U)\bigl(\alpha(v+U)+\beta(w+U)\bigr) &= (T/U)(\alpha v+\beta w+U)\\
    &= T(\alpha v+\beta w)+U\\
    &= \alpha Tv+\beta Tw+U\\
    &= \alpha(T/U)(v+U)+\beta(T/U)(w+U)
  \end{align*}
\item Suppose $\lambda$ is an eigenvalue of $T/U$ with corresponding
  eigenvector $v+U\in V/U$, so
  $(T/U)(v+U)=Tv+U=\lambda(v+U)=\lambda v+U$. This means $Tv-\lambda v\in U$.
  Suppose now that $T-\lambda I$ has no eigenvector in $V$ corresponding to
  $\lambda$. Since $U$ is invariant under $T-\lambda I$, then its restriction
  $(T-\lambda I)|_U$ is an operator on $U$. By the fundamental theorem of
  linear maps, we have
  \[ \dim U = \dim\nul(T-\lambda I)|_U+\dim\range(T-\lambda I)|_U. \]
  Now for any $v+U\ne U$ with $(T-\lambda I)v\in U$, then surjectivity of
  $(T-\lambda I)|_U$ would imply that there exists $u\in U$ such that
  $(T-\lambda I)u=(T-\lambda I)v$. By injectivity, we get that $u=v$, which is
  a contradiction since $v\notin U$. Therefore, our assumption that
  $T-\lambda I$ has no eigenvector corresponding to $\lambda$ must be false,
  and $\lambda$ is indeed an eigenvalue of $T$. \qedhere
\end{parts}
\end{solution}

\begin{exercise}
Suppose $V$ is finite-dimensional and $T\in\Lin(V)$. Prove that $T$ has an
eigenvalue if and only if there exists a subspace of $V$ of dimension
$\dim V-1$ that is invariant under $T$.
\end{exercise}
\begin{solution}
Suppose $T$ has an eigenvalue $\lambda$ with corresponding eigenvector
$v\in V$. Consider $\range(T-\lambda I)$. Then
$\dim\range(T-\lambda I)\le\dim V-1$. Extend to a subspace $U$ of $V$ of
dimension $\dim V-1$ that contains $\range(T-\lambda I)$. Then $U$ is invariant
under $T$ because for any $u\in U$, $Tu=(T-\lambda I)u+\lambda u\in U$. Hence,
there exists a subspace of $V$ of dimension $\dim V-1$ that is invariant under
$T$.

Conversely, suppose there exists a subspace $U$ of $V$ of dimension
$\dim V-1$ that is invariant under $T$. Then $V/U$ is one-dimensional, and the
induced operator $T/U$ on $V/U$ must have an eigenvalue $\lambda$. By
Exercise~5A38~(b), any eigenvalue of $T/U$ is also an eigenvalue of $T$. Hence,
$T$ has an eigenvalue.
\end{solution}

\begin{exercise}
Suppose $S,T\in\Lin(V)$ and $S$ is invertible. Suppose $p\in\Poly(\F)$ is a
polynomial. Prove that
\[ p(STS^{-1}) = Sp(T)S^{-1}. \]
\end{exercise}
\begin{solution}
We first prove by induction that $(STS^{-1})^k=ST^kS^{-1}$ for all
non-negative integers $k$. For $k=0$, we have $(STS^{-1})^0=I=ST^0S^{-1}$.
Suppose the result holds for some $k=n$, so that $(STS^{-1})^n=ST^nS^{-1}$.
Then
\begin{align*}
(STS^{-1})^{n+1} &= (STS^{-1})^n(STS^{-1})\\
  &= (ST^nS^{-1})(STS^{-1})\\
  &= ST^n(S^{-1}S)TS^{-1}\\
  &= ST^nTS^{-1}\\
  &= ST^{n+1}S^{-1}.
\end{align*}
By induction, the result holds for all non-negative integers $k$.

Now let $p(z)=a_0+a_1z+\cdots+a_mz^m$ be any polynomial. Then
\begin{align*}
p(STS^{-1}) &= a_0I+a_1(STS^{-1})+\cdots+a_m(STS^{-1})^m\\
  &= a_0I+a_1(STS^{-1})+\cdots+a_m(ST^mS^{-1})\\
  &= S(a_0I+a_1T+\cdots+a_mT^m)S^{-1}\\
  &= Sp(T)S^{-1}.
\end{align*}
Hence, the result holds for all polynomials $p$.
\end{solution}

\begin{exercise}
Suppose $T\in\Lin(V)$ and $U$ is a subspace of $V$ invariant under $T$. Prove
that $U$ is invariant under $p(T)$ for every polynomial $p\in\Poly(\F)$.
\end{exercise}
\begin{solution}
Let $u\in U$. Since $U$ is invariant under $T$, we have $T(u)\in U$. By
induction, suppose $T^k(u)\in U$ for some non-negative integer $k$. Then
$T^{k+1}(u)=T\bigl(T^k(u)\bigr)\in U$ since $U$ is invariant under $T$. Hence,
$T^k(u)\in U$ for all non-negative integers $k$. Now let
$p(z)=a_0+a_1z+\cdots+a_mz^m$ be any polynomial. Then
\[ p(T)(u) = a_0u+a_1T(u)+\cdots+a_mT^m(u) \in U, \]
since each term is in $U$. Therefore, $U$ is invariant under $p(T)$ for every
polynomial $p\in\Poly(\F)$.
\end{solution}

\begin{exercise}
Define $T\in\Lin(\F^n)$ by
$T(x_1,x_2,x_3,\dots,x_n)=(x_1,2x_2,3x_3,\dots,nx_n)$.
\begin{parts}
\item Find all eigenvalues and eigenvectors of $T$.
\item Find all subspaces of $\F^n$ that are invariant under $T$.
\end{parts}
\end{exercise}
\begin{solution}
\solnote{Appendix}{See the alternate proof at the end of this solution for a
generalization to any diagonal operator. The solution here remains as to show
the student the mechanism.}
\begin{parts}
\item Let $\lambda$ be an eigenvalue of $T$ with corresponding eigenvector
  $v=(x_1,x_2,\dots,x_n)\in\F^n$. Then $T(v)=\lambda v$, which means
  \[ (x_1,2x_2,3x_3,\dots,nx_n) = (\lambda x_1,\lambda x_2,\dots,\lambda x_n).
  \]
  Comparing components, we get $\lambda x_i=ix_i$ for each $i=1,2,\dots,n$.
  Hence, if $x_i\ne0$, then $\lambda=i$. Therefore, the eigenvalues of $T$ are
  $1,2,\dots,n$, and the corresponding eigenvectors are nonzero multiples of
  the standard basis vectors $e_1,e_2,\dots,e_n$.
\item Let $U$ be an invariant subspace of $\F^n$ under $T$. Let
  \[ u = \sum_{i=1}^n\alpha_ie_i \in U, \]
  where $e_i$ are the standard basis vectors of $\F^n$ and $\alpha_i\in\F$.
  Since $U$ is invariant under $T$, we have
  \[ Tu = \sum_{i=1}^n\alpha_iT(e_i) = \sum_{i=1}^n\alpha_iie_i \in U. \]
  Consider some index $i$ such that $\alpha_i\ne0$. Then
  \[ u_1 = Tu-iu = \sum_{j=1}^n\alpha_jje_j-i\sum_{j=1}^n\alpha_je_j
     = \sum_{j=1}^n\alpha_j(j-i)e_j \in U. \]
  The $i$-th term in this sum is $\alpha_i(i-i)e_i=0$, and $u_1$ now has all
  terms except the $i$-th term. Consider another index $j$ of $u_1$ such that
  $\alpha_j\ne0$. Then we can form
  \begin{align*}
  u_2 = Tu_1-ju_1
    &= \sum_{k=1}^n\alpha_k(k-i)ke_k-j\sum_{k=1}^n\alpha_k(k-i)e_k\\
    &= \sum_{k=1}^n\alpha_k(k-i)(k-j)e_k \in U.
  \end{align*}
  The $j$-th term in this sum is $\alpha_j(j-i)(j-j)e_j=0$, and $u_2$ now has
  all terms except the $i$-th and $j$-th terms. By repeating this process for
  each index with a nonzero coefficient, we can isolate each $\alpha_ie_i$ in
  $U$ whenever $\alpha_i\ne0$. Therefore, $U$ must be the span of some subset
  of the standard basis vectors $\{e_1,e_2,\dots,e_n\}$.
\end{parts}

\emph{Alternate proof.} \emph{Civil Service Pigeon.} Let $e_1,\dots,e_n$ be
the standard basis of $\F^n$.
\begin{parts}
\item For each $k\in\{1,\dots,n\}$, we have $Te_k=ke_k$. Hence, the eigenvalues
  of $T$ are $1,2,\dots,n$ and the eigenvectors corresponding to $k$ are the
  nonzero scalar multiples of $e_k$. Note that this list of eigenvalues (and
  eigenvectors) is exhaustive because $\dim(\F^n)=n$.
\item Let $U\subseteq\F^n$ be an invariant subspace under $T$. Take any
  $u=\sum_{i=1}^n\alpha_ie_i\in U$ and suppose $\alpha_k\ne0$ for some $k$. By
  Lagrange interpolation, there exists a polynomial $p$ such that
  $p(i)=\delta_{i,k}$ for all $i\in\{1,\dots,n\}$. Since $U$ is invariant under
  $T$, it is also invariant under $p(T)$, and so
  \[ p(T)u = \sum_{i=1}^n\alpha_ip(T)e_i = \sum_{i=1}^n\alpha_ip(i)e_i
     = \alpha_ke_k\in U \implies e_k\in U. \]
  Repeating this argument for every index with a nonzero coefficient in some
  element of $U$ shows that $U$ contains each of the basis vectors
  $e_1,\dots,e_n$. Hence, $U=\spn\{e_i : i\in S\}$ for some
  $S\subseteq\{1,\dots,n\}$.

  Conversely, any such span is invariant because $Te_i=ie_i$ belongs to that
  span. There are $2^n$ such subspaces, and they exhaust all invariant
  subspaces of $T$. \qedhere
\end{parts}
\end{solution}

\begin{exercise}
Suppose that $V$ is finite-dimensional, $\dim V>1$, and $T\in\Lin(V)$. Prove
that
$\{p(T) : p\in\Poly(\F)\}\ne\Lin(V)$.
\end{exercise}
\begin{solution}
Note that $\dim\Lin(V)\ge4$. Let $v_1,\dots,v_n$ be a basis for $V$. Consider
the operators $T_1,T_2\in\Lin(V)$ defined by
\[ T_1v_1=v_2,\quad T_1v_i=0\text{ for } i=2,\dots,n, \]
and
\[ T_2v_1=v_1,\quad T_2v_i=0\text{ for } i=2,\dots,n. \]
Then $T_2T_1v_1=T_2(T_1v_1)=T_2v_2=0$, and $T_1T_2v_1=T_1(T_2v_1)=T_1v_1=v_2$.
Hence, $T_2T_1\ne T_1T_2$, so $T_1$ and $T_2$ do not commute. Now let $S$ be
the set in question, and let $p,q\in\Poly(\F)$ where
\[ p(T)=a_0I+a_1T+\cdots+a_mT^m,\quad q(T)=b_0I+b_1T+\cdots+b_nT^n. \]
Observe that $p(T)q(T)$ consists of linear combinations of powers of $T$ with
coefficients $a_ib_j$. By commutativity of $\F$, we see that
$q(T)p(T)=p(T)q(T)$, so elements of $S$ commute. Since $T_1$ and $T_2$ do not
commute, they cannot both be in $S$. Hence, $S\ne\Lin(V)$.
\end{solution}
\end{exc}
